\section{Note taking Methods}
\label{research:NotetakingMethods}
\textbf{Source:} \cite{jesperSubjectGuidesNotetaking}

\subsection*{Outline Method}

\textbf{Method:}
\begin{itemize}
    \item Use a clear structure with high-level headings and indented subtopics and supporting facts.
    \item Subtopics should be indented to resemble code-style hierarchy.
    \item Works best with presentation notes, creating natural links to data.
\end{itemize}

\textbf{Pros and Cons:}

\begin{tabularx}{\textwidth}{>{\raggedright\arraybackslash}X>{\raggedright\arraybackslash}X}
    \rowcolor{gray!15} \textbf{Pros}                                                            & \textbf{Cons}                                                                  \\
    Ideal for structured subjects, making key concepts and supporting details easy to identify. & May encourage passive note-taking without critical reflection.                 \\
    Converts notes into study questions and improves active recall through indentation.         & Hard to insert new ideas without disrupting the structure.                     \\
    Great for digital notes; integrates well with online tools.                                 & Less suited to visual-heavy or problem-solving subjects like maths or physics. \\
    Enhances memory retention by clearly separating ideas from details.                         & Not ideal for fast-paced, unstructured lectures.                               \\
\end{tabularx}

\vspace{1em}
\subsection*{Cornell Method}

\textbf{Method:}
\begin{itemize}
    \item Created by Walter Pauk.
    \item Divide the page into three sections: cues, notes, and summary.
    \item Cue column references key points; notes contain details; summary reflects on content.
\end{itemize}

\textbf{Pros and Cons:}

\begin{tabularx}{\textwidth}{>{\raggedright\arraybackslash}X>{\raggedright\arraybackslash}X}
    \rowcolor{gray!15} \textbf{Pros}                                         & \textbf{Cons}                                            \\
    Encourages structured note-taking, improving organization and retrieval. & Not ideal for subjects with many diagrams (e.g., maths). \\
    Cue column supports self-testing and memory retention.                   & Space can feel cramped and hard to rearrange.            \\
    Summaries promote synthesis and deeper understanding.                    & Requires extra prep time before and after note-taking.   \\
    Structured format helps streamline review sessions.                      &                                                          \\
\end{tabularx}

\vspace{1em}
\subsection*{Mind Mapping Method}

\textbf{Method:}
\begin{itemize}
    \item Ideal for visual learners; mimics non-linear thinking.
    \item Begins at a central point and branches out radially.
    \item Allows ongoing addition of data to branches.
\end{itemize}

\textbf{Pros and Cons:}

\begin{tabularx}{\textwidth}{>{\raggedright\arraybackslash}X>{\raggedright\arraybackslash}X}
    \rowcolor{gray!15} \textbf{Pros}               & \textbf{Cons}                                               \\
    Highly visual overview of ideas.               & Not suited for capturing dense or detailed info.            \\
    Links ideas logically for easier navigation.   & Often more complex to set up; few good digital tools exist. \\
    Excellent for quick reviews and brainstorming. &                                                             \\
\end{tabularx}

\vspace{1em}
\subsection*{Charting Method}

\textbf{Method:}
\begin{itemize}
    \item Resembles a spreadsheet layout.
    \item Breaks information into categories and fills in columns.
    \item Needs well-defined structure upfront.
\end{itemize}

\textbf{Pros and Cons:}

\begin{tabularx}{\textwidth}{>{\raggedright\arraybackslash}X>{\raggedright\arraybackslash}X}
    \rowcolor{gray!15} \textbf{Pros}           & \textbf{Cons}                                              \\
    Great for minimal notes and comparisons.   & Text-based layout not ideal for visual subjects like STEM. \\
    Dense layout saves space and aids recall.  & Difficult to structure quickly during fast lectures.       \\
    Visual structure helps reinforce learning. & Time-consuming to set up effectively.                      \\
\end{tabularx}

\vspace{1em}
\subsection*{Sentence Method}

\textbf{Method:}
\begin{itemize}
    \item Involves writing shorthand, sequential notes.
    \item Captures the full lecture content but may lack structure.
    \item Requires later revision into a more structured format.
\end{itemize}

\textbf{Pros and Cons:}

\begin{tabularx}{\textwidth}{>{\raggedright\arraybackslash}X>{\raggedright\arraybackslash}X}
    \rowcolor{gray!15} \textbf{Pros}      & \textbf{Cons}                                        \\
    Consistent format for all subjects.   & Hard to differentiate primary vs. secondary ideas.   \\
    Easy to implement; no setup required. & Reviewing notes is inefficient—requires "unpacking". \\
    Compatible with digital tools.        & Can lead to note overload.                           \\
    No prep needed; works in real-time.   & May interfere with active listening.                 \\
\end{tabularx}

\vspace{1.5em}
\subsection*{My Takeaways}

After comparing multiple note-taking methods, it's clear that each has strengths and limitations. Subjects with distinct structures benefit from methods like Outline or Cornell, while creative or high-level understanding benefits from Mind Mapping. Personally, I find that combining \textbf{mind maps} for visual overviews with more structured formats like the \textbf{Cornell method} for deeper dives offers the best balance between flexibility and comprehension.

